% Fragment for inclusion in the main manuscript; no standalone preamble.
% This repairs the C<1 capacity case in the use of L23, Section 6.
\subsection{A two-factor repair for a capacity below one}
\label{sec:two-factor-greedy-repair}

In this subsection $\vartheta$ denotes the distribution exponent and
$\sigma$ the spectral parameter (called $\theta$ in L23).
The argument of L23, Lemma~6.3 and Proposition~6.5, uses a capacity
$C=D/(A b_0)$ after selecting a factor $b_0\le x^u$.
Lemma~6.3 itself assumes an existing factorization with $c\le C$;
its algebra is valid.  The subsequent greedy construction requires
care when $C<1$, because an integer factor $c\ge1$ cannot satisfy
$c\le C$.  Replacing this capacity by $\max(1,C)$ does not by itself
justify the three inequalities in L23~(6.3).  The following argument
handles that case by reallocating every prime factor of $d$.
It is the analogue, with the additional moving-residue constraints,
of the two-factor repair in Pascadi, arXiv:2505.00653v2,
Lemma~5.6 and the proof of Lemma~5.7.

\begin{lemma}[Two-factor repair]\label{lem:two-factor-greedy-repair}
Let $x>1$, $0<\delta<10^{-5}$, $1/2\le\vartheta\le5/8$,
$0\le\sigma\le7/32$, and $\alpha\ge0$.  Put
\[
 D=x^{\vartheta-2\delta},\qquad
 x^\delta\le A\le x^{1/3-\delta/2},\qquad B=D/A,
\]
and
\[
 u=\min\left\{
 \frac{2-(2+2\sigma)\vartheta}{3-4\sigma},
 \frac{2-2\vartheta-\alpha\sigma}{3}\right\}.
\]
Suppose $B\le x^u$ and
\[
 d=p_1\cdots p_r\in\mathcal D^{\rm well}(D),\qquad
 p_1\ge\cdots\ge p_r,
\]
where the defining condition is
$p_1\cdots p_{j-1}p_j^2\le D$ for every $j\le r$.
Then there is a factorization $d=a'b'c'$ with
$a'\le A$, $b'\le B$, $c'=1$, and
\begin{align}
 A^2b'c'^2&\le x^{1-3\delta},\label{eq:two-factor-repair-first}\\
 x^{\alpha\sigma}Aa'b'^5c'^2&\le x^{2-4\delta}
                       \le x^{2-3\delta},\label{eq:two-factor-repair-second}\\
 (Ac'/b')^{2\sigma}Aa'b'^5c'^2
   &\le x^{2-(4+4\sigma)\delta}
                       \le x^{2-2\delta}.
                       \label{eq:two-factor-repair-third}
\end{align}
These are the three constraints in L23~(6.3), with slack.
\end{lemma}

\begin{proof}
We have $A\ge1$ and
\[
 B\ge x^{\vartheta-1/3-3\delta/2}
       \ge x^{1/6-3\delta/2}>1.
\]
Start with $a'=b'=1$.  Immediately before assigning $p_j$, the
assigned product equals $a'b'=p_1\cdots p_{j-1}$.  The support
condition gives
\[
 p_j^2\le\frac{D}{a'b'}=\frac{A}{a'}\frac{B}{b'}.
\]
Consequently at least one of $p_j\le A/a'$ and $p_j\le B/b'$
holds.  Assign $p_j$ to that factor.  Iterating over all primes
preserves $a'\le A$, $b'\le B$ and yields $a'b'=d$.
This construction also covers $d=1$, by taking both factors equal
to one.  Set $c'=1$.

First,
\[
 A^2b'\le A^2B=AD
 \le x^{\vartheta+1/3-5\delta/2}
 \le x^{23/24-5\delta/2}
 \le x^{1-3\delta},
\]
where the last inequality follows from $\delta<10^{-5}<1/12$.
Next, using the second upper bound in the definition of $u$,
\[
 x^{\alpha\sigma}Aa'b'^5
 \le x^{\alpha\sigma}A^2B^5
 =x^{\alpha\sigma}D^2B^3
 \le x^{\alpha\sigma+2\vartheta-4\delta+3u}
 \le x^{2-4\delta}.
\]
Finally, $5-2\sigma>0$ and $3-4\sigma>0$, so the first upper
bound in the definition of $u$ gives
\begin{align*}
 (A/b')^{2\sigma}Aa'b'^5
 &\le A^{2+2\sigma}b'^{5-2\sigma}
 \le A^{2+2\sigma}B^{5-2\sigma}\\
 &=D^{2+2\sigma}B^{3-4\sigma}\\
 &\le x^{(2+2\sigma)\vartheta-(4+4\sigma)\delta
                         +(3-4\sigma)u}
 \le x^{2-(4+4\sigma)\delta}.
\end{align*}
\end{proof}

To apply the repair in the capacity case above, retain only the
information $b_0\le x^u$ and $C=D/(A b_0)<1$.  These imply
$B=D/A<b_0\le x^u$, which is the hypothesis of
Lemma~\ref{lem:two-factor-greedy-repair}.  The resulting $b'$ need
not equal $b_0$ or remain in $[x^v,x^u]$; the required conclusion is
the existence of some factorization satisfying L23~(6.3), already
proved directly.  If $C\ge1$, the original greedy construction
has two legitimate capacities and the usual Lemma~6.3 applies.
Thus this repair covers the missing case without reducing the
admissible distribution exponent, including the earlier level
$\vartheta=11/20$.

\paragraph{Correct initialization and capacities in the three-factor case.}
There is a separate issue in the printed proof of L23, Lemma~6.4:
the initialization must use $\{1,p_1,p_2\}$, and must append $p_3$
only in the subsequent greedy process.  Moreover, the printed
choice $D_2=D^2/x=x^{v-4\delta}$ does not equal $x^v$ and does not
give the claimed lower bound $b>x^v$ at the first failed step.
The capacities below avoid that discrepancy and give exact
versions of all three estimates.  We use the strict margins
already imposed on compact prefix boxes; no assertion at a gate
boundary without such margins is needed.

\begin{lemma}[Buffered three-factor construction]
\label{lem:buffered-three-factor-greedy-repair}
Use the notation and ranges of
Lemma~\ref{lem:two-factor-greedy-repair}, without assuming
$B\le x^u$, and let $v=2\vartheta-1$.
Let $d=p_1\cdots p_r\in\mathcal D^{\rm well}(D)$ with
$p_1\ge\cdots\ge p_r$.
Suppose $r\ge3$, $p_3\le x^{u-v}$, and, for some
$\{i,j\}=\{1,2\}$,
\[
 p_i\le x^v,\qquad p_j^2\le x^{1-\vartheta-2\delta}.
\]
Then $d$ admits a factorization satisfying L23~(6.3).
\end{lemma}

\begin{proof}
Set
\[
 D_1=A,\qquad D_2=x^v,\qquad D_3=\frac{D}{Ax^v}.
\]
Their product is exactly $D$.  All three capacities are at least
one: for the last one,
\[
 D_3\ge x^{2/3-\vartheta-3\delta/2}
       \ge x^{1/24-3\delta/2}>1.
\]
Place $p_i$ in the second factor.  Since
$p_j^2\le D_1D_3$, place $p_j$ in either the first or third
factor that has sufficient capacity, leaving the other equal to
one.  The assigned factors are thus a permutation of
$\{1,p_1,p_2\}$ with the stated placement of $p_i$.
Now append $p_3,p_4,\ldots$ whenever any factor has capacity.

If every prime is appended, write the resulting factors as
$a,b,c$ with $a\le D_1$, $b\le D_2$, $c\le D_3$.
Since $p_3>1$ and $p_3\le x^{u-v}$, we have $u\ge v$.
The two inequalities defining $u\ge v$ imply, respectively,
\[
 (8-6\sigma)\vartheta+4\sigma\le5,
 \qquad \alpha\sigma+8\vartheta\le5.
\]
Consequently
\begin{align*}
 A^2bc^2
 &\le A^2D_2D_3^2
   =x^{1-4\delta},\\
 x^{\alpha\sigma}Aab^5c^2
 &\le x^{\alpha\sigma}A^2D_2^5D_3^2
   =x^{\alpha\sigma+8\vartheta-3-4\delta}
   \le x^{2-4\delta},\\
 (Ac/b)^{2\sigma}Aab^5c^2
 &\le A^{2+2\sigma}D_2^{5-2\sigma}D_3^{2+2\sigma}\\
 &=x^{(8-6\sigma)\vartheta+4\sigma-3-(4+4\sigma)\delta}
   \le x^{2-(4+4\sigma)\delta}.
\end{align*}
These imply the three required inequalities.

Otherwise let $p_k$, with $k\ge3$, be the first prime that cannot
be appended, and denote the current factors by $a_0,b_0,c_0$.
Thus $a_0b_0c_0=p_1\cdots p_{k-1}$ and
$a_0p_k>A$, $b_0p_k>x^v$, $c_0p_k>D_3$.
Put $b=b_0p_k$.  Then
\[
 x^v<b\le x^v p_3\le x^u.
\]
The support condition and $a_0p_k>A$ give
\[
 A b c_0<(a_0p_k)(b_0p_k)c_0\le D.
\]
It follows that $C=D/(Ab)>c_0\ge1$, so both remaining
capacities $A,C$ are legitimate.  For every remaining prime
$p_l$ the support condition gives
\[
 p_l^2\le
 \frac{D}{a_0bc_0p_{k+1}\cdots p_{l-1}}
 =\frac{AC}{a_0c_0p_{k+1}\cdots p_{l-1}}.
\]
The two-factor greedy argument, starting from $a_0,c_0$, therefore
assigns all remaining primes and produces $d=abc$ with
$a\le A$, $c\le C=D/(Ab)$, and $x^v<b\le x^u$.
Directly,
\begin{align*}
 A^2bc^2&\le D^2/b\le x^{1-4\delta},\\
 x^{\alpha\sigma}Aab^5c^2
 &\le x^{\alpha\sigma}D^2b^3\le x^{2-4\delta},\\
 (Ac/b)^{2\sigma}Aab^5c^2
 &\le D^{2+2\sigma}b^{3-4\sigma}
   \le x^{2-(4+4\sigma)\delta}.
\end{align*}
This also covers a failed step at $k=r$, when no primes remain
to be reassigned.
\end{proof}

For a strict $W$ branch, its inequalities state that the
designated prefix exponent is below $v$, that the other prefix
prime has square below $x^{1-\vartheta}$, and that
$t_3<u-v$.  On a fixed compact box in this strict region, choose
$\delta$ small enough that the second inequality still holds
with exponent $1-\vartheta-2\delta$; shrink the other margins to
cover the box widths and the fixed perturbation of $\alpha$.
Lemma~\ref{lem:buffered-three-factor-greedy-repair} then applies.
As these positive buffers tend to zero, the limiting eight-branch
envelope $\Phi$ is unchanged.  This is a buffered use of the
source envelope, not a claim that its printed boundary argument
or its printed initialization was exact.
